%% ============================================================
%%  SECCIÓN 2 — Resumen
%% ============================================================

\section{Resumen}
\label{sec:resumen}

En este informe presentamos la caracterización experimental y el análisis cuantitativo de los procesos transitorios de carga y descarga en circuitos RC de primer orden, desarrollada en la Mesa C7 del laboratorio de Física III de la Universidad Nacional de Ingeniería. El montaje experimental se alimentó con una señal de onda cuadrada periódica simétrica de amplitud \qty{10,0}{\voltpp} y frecuencia \qty{250}{\hertz} suministrada por un generador de funciones arbitrario OWON AG 051F, asegurando que el semiperíodo de excitación ($T/2 = \qty{2,00}{\milli\second}$) fuera holgadamente superior a cinco constantes de tiempo ($T/2 \ge 5{,}8\tau$) para permitir la saturación asintótica completa del condensador.

Utilizando un osciloscopio digital Rohde \& Schwarz HMO1202 con base de tiempo calibrada en \qty{500}{\micro\second\per\division}, medimos de forma simultánea la tensión en el condensador $V_C(t)$ y la caída de potencial en la resistencia $V_R(t)$ para combinaciones sistemáticas de tres resistores patronados ($R_1 = \qty{9,980 \pm 0,058}{\kilo\ohm}$, $R_2 = \qty{6,832 \pm 0,033}{\kilo\ohm}$ y $R_3 = \qty{2,250 \pm 0,012}{\kilo\ohm}$) y tres condensadores comerciales ($C_1, C_2, C_3$). Determinamos la constante de tiempo experimental $\tau$ en la cota del \qty{63,2}{\percent} de la tensión de régimen permanente. 

A través de regresiones lineales por mínimos cuadrados ordinarios (MCO) sobre el modelo teórico $\tau = m R + b$, obtuvimos las siguientes capacitancias experimentales:
\begin{align*}
    C_{1,\text{OLS}} &= \qty{10,79 \pm 0,05}{\nano\farad} \quad (R^2 = 0{,}99998; \text{ error rel. } \qty{+7,95}{\percent}), \\
    C_{2,\text{OLS}} &= \qty{27,44 \pm 0,03}{\nano\farad} \quad (R^2 = 0{,}99999; \text{ error rel. } \qty{+22,48}{\percent}), \\
    C_{3,\text{OLS}} &= \qty{34,61 \pm 0,29}{\nano\farad} \quad (R^2 = 0{,}99993; \text{ error rel. } \qty{+4,15}{\percent}).
\end{align*}
Al representar $\tau$ frente a $R$, se constató una rigurosa proporcionalidad directa ($\tau \propto R$) donde las rectas pasan prácticamente por el origen de coordenadas: los interceptos residuales se mantuvieron en fracciones de microsegundo ($|b| < \qty{0,3}{\micro\second}$), plenamente compatibles con cero dentro de la resolución de los cursores (\qty{0,1}{\micro\second}) y la incertidumbre instrumental del osciloscopio HMO1202, corroborando la ley física fundamental $\tau = RC$. La mayor discrepancia observada en $C_2$ (\qty{+22,48}{\percent}) se atribuye a la tolerancia comercial de fabricación ($\pm\qty{20}{\percent}$) y a la degradación por envejecimiento dieléctrico.

Asimismo, mediante la función aritmética \texttt{MATH ADD} del osciloscopio digital comprobamos en tiempo real la Segunda Ley de Kirchhoff ($V_C(t) + V_R(t) = V_{\text{gen}}(t)$), registrando una onda escalón perfectamente plana con una fluctuación residual inferior a \qty{0,05}{\volt}. Finalmente, para la red mixta con divisor resistivo del paso 20, verificamos el Teorema de Thévenin al corroborar que la tensión máxima en el capacitor se reduce a $V_{\text{th}} = \qty{4,06}{\volt}$ y que la corriente suministrada decae exponencialmente hacia un valor asintótico no nulo de \qty{0,59}{\milli\ampere}.
